Мета роботи: дослідження принципів формування криптографічно стійких випадкових 
послідовностей, генерації ключового матеріалу, тестування та генерації простих чисел 
для асиметричних криптосистем, а також основних вимог до життєвого циклу 
криптографічних ключів. 

\section{Дослідження принципів генерації криптографічних ключів}

\subsection{Що таке ентропія у контексті генерації криптографічних ключів? Які джерела ентропії можуть використовувати сучасні програмно-апаратні системи?}

Ентропія у контексті криптографії характеризує ступінь невизначеності випадкових даних, які використовуються для створення криптографічних ключів. 
Чим вища ентропія, тим складніше зловмиснику передбачити значення ключа або підібрати його перебором. Навіть якщо ключ має достатню довжину, його безпека може бути низькою, 
якщо він згенерований на основі передбачуваних даних.

Сучасні програмно-апаратні системи можуть використовувати різні джерела ентропії. До них належать фізичні процеси, зокрема тепловий шум електронних компонентів 
і нестабільність частоти коливань електронних схем. Також можуть враховуватися часові характеристики системних подій, рухи миші, натискання клавіш та інші непередбачувані події. 
Операційні системи можуть поєднувати дані з кількох джерел, обробляти їх і використовувати для ініціалізації криптографічно стійких генераторів випадкових чисел. 
При цьому важливо не лише збирати випадкові дані, а й оцінювати їхню якість та рівень ентропії \cite{RFC4086,NIST80090B}.

\subsection{У чому полягає різниця між TRNG, PRNG та криптографічно стійким генератором псевдовипадкових чисел (CSPRNG)?}

TRNG (True Random Number Generator) є генератором справжніх випадкових чисел, який використовує фізичні процеси для отримання випадкових значень. 
Такі генератори можуть спиратися на тепловий шум або інші фізичні явища, результати яких важко передбачити. Їхньою особливістю є те, що випадковість формується 
на основі фізичного джерела, а не лише обчислень.

PRNG (Pseudorandom Number Generator) є генератором псевдовипадкових чисел, який за певним алгоритмом формує послідовність значень на основі початкового стану. 
Такі послідовності можуть мати хороші статистичні властивості, однак за однакових початкових умов генератор відтворює ті самі значення.

CSPRNG (Cryptographically Secure Pseudorandom Number Generator) є криптографічно стійким генератором псевдовипадкових чисел. Він також використовує алгоритм 
для формування послідовності, але має додаткову вимогу щодо непередбачуваності результатів. Знання попередніх значень не повинно давати змоги ефективно визначити 
наступні значення або відновити внутрішній стан генератора за умови дотримання вимог до його реалізації та ініціалізації. 
CSPRNG зазвичай використовує достатньо випадкове початкове значення, отримане з надійного джерела ентропії. Саме такі генератори застосовують для створення криптографічних 
ключів та інших секретних даних \cite{RFC4086,NIST80090A}.

\subsection{Яку роль виконує початкове значення (seed) генератора псевдовипадкових чисел? Які наслідки може мати можливість передбачити або відтворити seed?}

Початкове значення (seed) використовується для ініціалізації генератора псевдовипадкових чисел. На його основі формується початковий стан, 
з якого алгоритм надалі генерує послідовність значень. Для звичайного детермінованого генератора однакове початкове значення за однакових умов призводить 
до відтворення тієї самої послідовності. Це є корисним під час моделювання та проведення експериментів, оскільки дозволяє повторити отримані результати.

У криптографії можливість передбачити або відтворити seed становить серйозну загрозу безпеці. Якщо зловмисник знає початкове значення або може перебрати невелику 
кількість можливих варіантів, він може відтворити послідовність генератора та визначити створені на її основі ключі. У такому випадку навіть ключ великої довжини може виявитися вразливим, 
оскільки фактична кількість можливих ключів обмежується невизначеністю початкового значення. Тому для криптографічних генераторів важливо використовувати достатньо випадковий seed, 
захищати внутрішній стан і правильно реалізовувати механізми ініціалізації та оновлення стану \cite{RFC4086,NIST80090A}.

\subsection{Чому звичайні генератори псевдовипадкових чисел, призначені для моделювання, ігор або статистичних задач, не повинні використовуватися для генерації криптографічних ключів?}

Звичайні генератори псевдовипадкових чисел призначені насамперед для отримання послідовностей, які мають необхідні статистичні властивості та підходять для моделювання, 
ігор або статистичних експериментів. Проте такі генератори не обов'язково забезпечують захист від спроб передбачити наступні значення або відновити їх за відомими результатами. 
Тому статистично випадкова на вигляд послідовність не обов'язково є безпечною для використання в криптографії.

Якщо криптографічний ключ створено за допомогою звичайного генератора, зловмисник може скористатися особливостями алгоритму, передбачуваним початковим значенням 
або відомими результатами його роботи, щоб звузити множину можливих ключів. Це може призвести до розкриття зашифрованих даних або компрометації системи автентифікації. 
Наприклад, у мові Python модуль \texttt{random} призначений для моделювання та статистичних задач, тоді як модуль \texttt{secrets} надає засоби для отримання криптографічно 
стійких випадкових значень. Отже, для генерації криптографічних ключів потрібно використовувати спеціально призначені криптографічно стійкі генератори, які спираються 
на належне джерело випадковості та відповідають вимогам безпеки \cite{RFC4086,PythonRandom,PythonSecrets}.


\subsection{Яким чином сучасні операційні системи надають програмам криптографічно стійкі випадкові значення? Навести приклад відповідного API для однієї операційної системи або мови програмування.}

\subsection{Чим принципово відрізняється генерація випадкового симетричного ключа від генерації пари ключів асиметричної криптосистеми?}

\subsection{Чи гарантує використання алгоритму AES-256 рівень захисту у 256 біт, якщо ключ отримано з недостатньо випадкового або передбачуваного джерела? Обґрунтувати відповідь.}

\subsection{Які основні етапи життєвого циклу криптографічного ключа необхідно враховувати, крім безпосередньої генерації ключа?}
